\section{Caso de ingeniería: repositorios reales con participación de Agent}
\subsection{Commits semiautomatizados}
Graham compartió su práctica con los repositorios: en algunos proyectos, «la mitad del contenido de los commits lo redacta primero un Agent de codificación, y después él lo revisa y corrige». El Agent primero lee el Issue, encuentra los archivos relevantes mediante resumen y búsqueda, después redacta un parche preliminar y, por último, un ingeniero humano lo integra.

Él prevé que, en el próximo año o dos, los Agent se encargarán de más tareas cotidianas y rutinarias, para que los ingenieros puedan concentrar su atención en el diseño de alto nivel. Por ahora, el Agent puede manejar sin problemas los scripts de Package con dependencias compartidas, las operaciones básicas de CLI y las correcciones simples de errores (bug fix); los escenarios complejos todavía requieren revisión humana.

\subsection{Revisión del proceso y retroalimentación}
Después de cada ejecución del Agent, el equipo guarda un registro final: incluye los archivos obtenidos mediante fetch de forma proactiva, los comandos invocados, los resultados de las pruebas y las causas de los fracasos. Estos registros se convierten en muestras de entrenamiento y recordatorios de retrollamada, para que el Agent aprenda a incorporar en la cadena de prompts los pasos útiles de las muestras fallidas.

También recalcó: «Observar el proceso de ejecución del Agent y dar correcciones inmediatas mejora más la calidad del modelo que puntuar únicamente el resultado.» Esta retroalimentación basada en la observación incluye pausas humanas, la inyección de prompts y hacer que el Agent vuelva a intentarlo después de corregirlo.

\begin{knowledgebox}{Experiencia del caso}
Graham ya usa el Agent para redactar en varios repositorios, y los commits que aporta el Agent representan cerca de la mitad del volumen de trabajo. Esto demuestra que el Agent ya tiene un valor práctico en las tareas repetitivas.
\end{knowledgebox}

\begin{importantbox}{Recomendación práctica}
Despliegue primero al Agent en tareas de menos de 10 minutos (como el formateo, los parches o las actualizaciones de dependencias), arme un proceso de review humano y, después, extiéndalo poco a poco hacia longer-running job.
\end{importantbox}

\begin{warningbox}{No deje que el Agent se despliegue solo}
Aun cuando el Agent pueda completar ciertos cambios, no se le debe dejar llevarlos solo a producción: deben conservarse la revisión humana, las pruebas de regresión y la aprobación (approve) manual.
\end{warningbox}

\subsection{Resumen de la sección}
\begin{itemize}
    \item Los experimentos con el Agent en repositorios reales ya muestran el potencial de la semiautomatización, pero todavía no bastan para sustituir por completo a los ingenieros humanos.
    \item Llevar un registro del proceso y corregir de inmediato son las prácticas fundamentales para mejorar la confiabilidad del Agent.
\end{itemize}
